\documentclass[11pt]{article}
\usepackage{amsmath,hyperref}
\title{A Moving Inspection Interval Cannot Cross a Gap at a Phase Knot}
\date{2026-11-27}
\begin{document}
\maketitle

\paragraph{Pins.}
This is a deterministic simulator-level reference for known
piecewise-affine requirements in a single-file corridor. It uses the
previous closed difference-constraint sensing, radio and clearance
cells, with order, speed bounds and health fixed over a phase. Time
windows and motion of the protected interval are prescribed. The
runtime does not learn the moving window or fault identity here.
Bodies and radio are modeled by the given inequalities; continuous
coverage of a modeled interval is not telemetry quality or hardware
delivery. No acceleration or rest constraint is imposed.

\paragraph{Time lift.}
Use coordinates \((t,x)\), where \(t\) is the physical clock, not a
free curve parameter. On each closed interval of positive duration
write the inspection-window shift as \(s(t)=u t+d\). Add \(s(t)\)
to its upper endpoint bounds, subtract it from lower endpoint bounds,
and leave absolute arena, separation and radio inequalities in
their original coordinates. Every old sensing cell becomes a closed
convex polyhedron in \(t,x\). The accepted model is their union
within the prescribed phases. Where a phase boundary is shared,
continuity and closedness force the two one-sided limits to
satisfy their corresponding requirements. A separate mandatory
zero-duration event would need its own explicit cut, not a
zero-width cell treated as optional.

The last-visit proof of the preceding continuous static result
applies unchanged to the lifted path \((t,x(t))\). Choose last
visits using physical time. Chords have that same elapsed time
between endpoints, since the first coordinate interpolates from
\(t_j\) to \(t_{j+1}\) at unit rate; a zero-time chord has equal
position by the speed inequalities. Thus exact feasibility at a
deadline reduces to finitely many simple lifted cell sequences with
nondecreasing time and
\[
 |x_{j+1,i}-x_{j,i}|\le v_i(t_{j+1}-t_j).
\]
This does not hand the controller a post-fault truth before its
observations: the program tests an exogenous known timed plant.
The held-in exact certificate harness checks rational feasible
points and rational Farkas witnesses, refusing reconstruction
failure rather than rounding a verdict. Independent chord coverage
subtracts the known shift on each phase and tests the intervals
in relative coordinates. General exact rational linear programming
can replace the numerical proposal step.

\paragraph{Sharp two-robot discriminator.}
Take robots at \(0,4\) at both \(t=0\) and \(t=4\).
The first is stationary, the second has speed at most one, and their
gap must stay in \([1,6]\); sensing radius is two and ordinary single
coverage is required at every point of \([s(t),s(t)+1]\).
For peak \(h\geq0\), let
\(s(t)=h\min(t,4-t)/2\). The phase knot \(t=2\) alone allows the
moving robot at most \(6\) and hence a covered far edge at most
\(8\); knots permit exactly \(h\le7\). In particular a physically
speed-feasible trajectory through those three placements is not
yet covered throughout.

For \(h>2\), while \(1<s(t)<2\) on the outward leg the interval
straddles the end \(2\) of the fixed footprint. The movable
robot cannot leave a gap at that seam, so its position is at most
\(4\). At \(t_c=4/h\), continuity still enforces \(x_2(t_c)\le4\).
Between this cut and the peak the robot can advance by at most
\(2-4/h\). At the peak its far sensing edge must cover \(h+1\),
hence
\[
 h+\frac4h\le 7.
\]
Conversely, hold at \(4\) through the seam and thereafter track
\[
 x_2(t)=4+\max\{0,\min(t,4-t)-4/h\}.
\]
This speed-feasible trajectory has gap in \([4,6]\).
Before the crossing the two footprints touch and cover the whole
required interval. Between \(t_c\) and the peak the left required
endpoint moves at slope \(h/2>1\) and never falls behind the
moving footprint's left endpoint. Its right endpoint remains
covered because the tightest right-hand inequality is at the
peak. The return is symmetric. If \(h\le2\), static positions
are sufficient. Thus the exact maximum is
\((7+\sqrt{33})/2\), not 7; Lean checks the scalar polynomial
equivalence in the relevant regime. At \(h=6.5\), the three
samples at \(t=0,2,4\) can all pass but no continuous coverage
trajectory exists, even with the entire window known. This is
physical infeasibility in the model, not an information delay.

\paragraph{Evidence and limits.}
Across 25 rational peaks and 307 lifted routes the feasibility
decoder agrees with the analytic inequality, and constructive
positive paths pass exact relative-interval chord checks.
The bound is input to a future scheduling or control proof, not a
guarantee that a mobile crew sees the shift on time. A direct
spacetime linear planner supplied the same union of cells can
perform the same exact check. No novelty or contract acceptance
is claimed over that mechanism, no independent mission is used,
and compatible continuations under multiple unobserved future
requirements remain to be tested.
\end{document}
